% =========================================================================
% Capítulo 1 - Tema 01: Cinemática en Coordenadas Cilíndricas
% Problema 04: Disco Rotante con Punto en Movimiento Radial
% =========================================================================
% Archivo: C1_E01_P04_disco_rotante.tex
% Descripción: Movimiento de un punto que se desplaza radialmente sobre
%              un disco que gira a velocidad angular constante
% =========================================================================

\subsubsection{Disco Rotante con Punto en Movimiento Radial}

Un disco de radio $R$ gira con velocidad angular constante $\Omega$ respecto a un marco de referencia fijo. Un punto $P$ se mueve radialmente sobre el disco de acuerdo con la ley $r(t) = r_0 + vt$, donde $r_0$ es la posición radial inicial y $v$ es la velocidad radial constante. El punto siempre permanece sobre el disco (sin separarse). Observe el movimiento desde el marco de referencia fijo.

\begin{enumerate}
    \item Determinar la velocidad del punto $P$ en coordenadas cilíndricas.
    \item Encontrar la aceleración del punto $P$ en coordenadas cilíndricas.
    \item Identificar las componentes de aceleración: radial (centrípeta), transversal y axial.
\end{enumerate}

\begin{figure}[!ht]
    \centering
    \begin{tikzpicture}[>=Latex, scale=1.8]

        % --- CONFIGURACIÓN DE PARÁMETROS ---
        \def\R{1.8}      % Radio del disco [cm]
        \def\r{1.0}      % Posición actual del punto [cm]
        \def\OmegaAng{35} % Ángulo de rotación actual [grados]

        % --- SOPORTES Y ELEMENTOS FIJOS ---
        % Eje de rotación (estacionario)
        \draw[thick] (-0.3, -2.5) -- (-0.3, 0.5) node[left, font=\small] {Eje fijo};

        % --- GEOMETRÍA PRINCIPAL ---
        % Disco (círculo)
        \draw[thick, double distance=2pt] (0,0) circle (\R);
        \coordinate (O) at (0,0);
        \filldraw (O) circle (1.5pt) node[left] {$O$};

        % --- SISTEMAS DE COORDENADAS ---
        % Marco fijo (ejes x-y)
        \draw[->, thick, black] (O) -- (2, 0) node[right, font=\small] {$x$};
        \draw[->, thick, black] (O) -- (0, 2) node[above, font=\small] {$y$};

        % --- PUNTOS Y COORDENADAS ---
        % Posición del punto P
        \coordinate (P) at ({\OmegaAng}:\r);
        \filldraw[fill=white, thick] (P) circle (2.5pt);
        \filldraw (P) circle (1pt) node[above right] {$P$};

        % Línea radial (radio instantáneo)
        \draw[dashed, gray] (O) -- ({\OmegaAng}:1.9);

        % --- VECTORES Y FUERZAS ---
        % Velocidad (componente radial)
        \draw[->, thick, azul] (P) -- ($(P) + ({0.6*cos(\OmegaAng)}, {0.6*sin(\OmegaAng)})$)
            node[right, font=\small, color=azul] {$\vec{v}_r$};

        % Velocidad (componente transversal - tangencial)
        \draw[->, thick, myred] (P) -- ($(P) + ({-0.6*sin(\OmegaAng)}, {0.6*cos(\OmegaAng)})$)
            node[above, font=\small, color=myred] {$\vec{v}_\theta$};

        % --- ANOTACIONES Y ETIQUETAS ---
        % Ángulo theta
        \draw[->, dashed] (0.4, 0) arc (0:\OmegaAng:0.4);
        \node at ({\OmegaAng/2}:0.65) {$\theta = \Omega t$};

        % Distancia radial
        \draw[<->, thin] (0, -0.3) -- ({\OmegaAng}:\r) node[midway, below left, font=\small] {$r(t)$};

        % Ángulo de rotación
        \node[font=\footnotesize] at (0.3, -2) {$\omega = \Omega = \text{const}$};

    \end{tikzpicture}
    \caption{Disco rotante con punto que se desplaza radialmente; velocidad descompuesta en componentes cilíndricas}
    \label{fig:c1_e01_p04_disco_rotante}
\end{figure}

\newpage

\subsubsection*{Notas sobre este problema}

Este ejercicio ilustra un caso especial donde el movimiento puede parecer sencillo en coordenadas cartesianas (el disco está girando uniformemente), pero revela su complejidad cuando se analiza desde coordenadas cilíndricas. El punto tiene una velocidad radial constante $v$ (que el observador fijo ve muy claramente), pero también experimenta una velocidad tangencial debido a la rotación del disco. La clave es entender cómo estas dos velocidades se combinan para producir la aceleración total.
